\section[A homogeneous scalar lift at the unbounded-state contact]{A homogeneous scalar lift at the unbounded-state contact}
\label{sec:unbounded-b-homogeneous-scalar}

This lemma is a scalar lifting statement.  It neither supplies the full
state normals nor proves positivity of its coefficient Gram by itself.
Its purpose is to avoid any bound on the contact state ratio.

Let $B,\sigma>0$ be pure coefficient infinitesimals, let $k>0$ be a
pure coefficient bounded unit, and put $x=\sigma B$, $\xi=u^2$.
Suppose
\[
 f=a(\xi,v)W^2+2u d(\xi,v)WZ+c(\xi,v)Z^2,
 \quad
 a=\frac\xi B Q+\sigma^2 a_0,\qquad
 c=B e+\xi b,                                      \tag{HL1}
\]
where $Q,a_0,d,e,b$ have bounded polynomial coefficients.  Bounded
pure testing unit denominators may be cleared beforehand.
Assume the value and all first derivatives of $f$ vanish at
\[
 u=\sqrt{x},\quad v=k,\quad
 W=\frac1k\sqrt{B/\sigma},\quad Z=1.                \tag{HL2}
\]
No bound on $B/\sigma$ is assumed.

Define the homogeneous roots
\[
 R=\frac{uW-(B/k)Z}{\sqrt B},\qquad
 S=uZ-\sigma kW,\qquad
 U=\sqrt B(v-k)Z,\qquad T=u(v-k)Z.                 \tag{HL3}
\]
They satisfy the exact conductors
\[
 (\xi-x)W=u\sqrt B R+(B/k)S,
 \qquad
 (\xi-x)Z=uS+\sigma k\sqrt B R.                    \tag{HL4}
\]
Then $f$ has an exact symmetric Gram on $(R,S,U,T)$ whose entries
are bounded polynomials.  In particular no inverse of $\sigma$,
$B$, a testing infinitesimal, or a mixed expression is required in
these Gram entries.  The roots themselves need not be bounded.

\paragraph{First set $v=k$.}
In this paragraph every scalar coefficient is evaluated at $v=k$.
State stationarity and $u$ stationarity give
\begin{gather*}
 a_c=-\sigma k d_c,\qquad c_c=-B d_c/k,\\
 B(a_\xi)_c+kd_c+2kx(d_\xi)_c
                      +\sigma k^2(c_\xi)_c=0,     \tag{HL5}
\end{gather*}
where $c$ on a subscript means evaluation at $\xi=x$.
These are identities in the pure coefficient field; deriving them
does not insert an inverse into a tensor Gram.

Set, by monic polynomial division in $\xi$,
\[
 s=\frac{c+Bd/k}{\xi-x},\qquad
 q=-kd-\sigma k^2s,\qquad
 h_2=\frac{Ba+k\xi d+\sigma k^2c+\sigma Bkd}
                    {(\xi-x)^2}.                 \tag{HL6}
\]
All three quotients have bounded polynomial coefficients.  The first
numerator vanishes at $x$ by (HL5), and the last numerator and its
$\xi$ derivative vanish there by the same three identities.
Moreover $Ba=\xi Q+B\sigma^2a_0$ is bounded.  Division by a monic
polynomial preserves bounded coefficients and is not evaluation of
an inverse of $\xi-x$.

Direct multiplication gives the exact identity
\[
 f(\xi,k)=qR^2+sS^2
                  +h_2\left(uR+\frac{\sqrt B}{k}S\right)^2.
                                                               \tag{HL7}
\]
For later residue computations, if $E_0=e+d/k$, then
\[
 s=b+B\frac{E_0(\xi)-E_0(x)}{\xi-x}
       +x\frac{b(\xi)-b(x)}{\xi-x}.                \tag{HL8}
\]
Thus $s_0=b_0$, $q_0=-k_0d_0$, and the $h_2$ correction has zero
residue in every Gram entry.  Here subscripts $0$ denote residues,
not evaluation at the contact polynomial variable.

\paragraph{The complete angular coquotient.}
Write $\tau=v-k$ and, for each of $Q,a_0,d,e,b$, form its bounded
monic difference in $v$, denoted respectively by
$Q_1,a_{01},d_1,e_1,b_1$:
\[
 Q(\xi,v)-Q(\xi,k)=\tau Q_1(\xi,v),
\]
and similarly for the other four functions.  Put
\[
 H_1=Q_1/k^2+2d_1/k+e_1,\qquad
 L_1=a_{01}/k^2+b_1,\qquad J_1=H_1+\sigma L_1.
\]
Angular stationarity at (HL2) gives exactly $J_1(x,k)=0$.
Consequently the successive monic differences
\[
 J_1=(\xi-x)J_\xi+\tau J_v,
 \quad
 J_\xi=\frac{J_1(\xi,v)-J_1(x,v)}{\xi-x},\qquad
 J_v=\frac{J_1(x,v)}{v-k}                            \tag{HL9}
\]
are bounded polynomials.  Let $D_1=L_1+BJ_\xi$.  The full exact
angular correction is
\begin{align*}
 f(\xi,v)-f(\xi,k)={}&
 \tau Q_1R^2+\frac{\tau a_{01}}{k^2}S^2\\
 &+\left\{2(Q_1/k+d_1)+\sigma kD_1\right\}RU\\
 &+\left\{D_1-2a_{01}/k^2\right\}ST+J_vU^2.       \tag{HL10}
\end{align*}
To check this without any division, use
$uW/\sqrt B=R+\sqrt B Z/k$ and
$\sigma W=(uZ-S)/k$.  The only term not already among the displayed
root products becomes
\[
 \tau Z^2\{BH_1+\xi L_1\}
 =D_1\tau(\xi-x)Z^2+BJ_v\tau^2Z^2
 =D_1(ST+\sigma kRU)+J_vU^2,
\]
by (HL4).  Equations (HL7) and (HL10) prove the asserted homogeneous
Gram lift for every state $W,Z$ simultaneously.  There is no state
normalization and no discarded endpoint $Z=0$.

The usual exact zero, now with the new source scale, is
\[
 T^2-\sigma U^2-ST(v-k)-\sigma kRU(v-k)=0.          \tag{HL11}
\]
This is the appropriate identity for redistributing the angular
diagonal.  Its $U^2$ coefficient is $\sigma$, which is infinitesimal;
replacing it by $1$ would erase the smaller radial scale.  Any
positivity argument must use the actual residue coefficients of
(HL7)--(HL10), rather than importing the bounded-state Schur test
without this scale correction.

\paragraph{Strict scalar corollary for the small-$B$ source.}
Suppose additionally that the coefficientwise leading functions on
the finite $u=0$ axis are
\[
 Q_0(v)=\frac{2K^2(1+v^2)}{1+K^2},\qquad
 d_0(v)=-2v,\qquad
 e_0(v)=\frac{2(v^2/K^2+K^2)}{1+K^2},
 \qquad b_0(K)>0,
\]
where $K=k_0$.  Then the $(R,U)$ residue block produced above is
\[
 2K^2R^2+\frac{4(K^2-1)}{1+K^2}RU+\frac2{K^2}U^2,
\]
with Schur margin $8/(1+K^2)^2>0$.  The $S^2$ coefficient has
residue $b_0(K)>0$, the $RS$ residue is zero, and the only other
possible residue entry is a bounded $ST$ coefficient, say $\ell$.
Add a positive real multiple $\delta$ of (HL11), with
$\delta>\ell^2/(4b_0(K))$.  Its correction to the $U^2$ entry has
zero residue because $\sigma$ is infinitesimal, whereas it supplies
the positive $T^2$ diagonal $\delta$.  Both residue blocks are then
strictly positive.  This proves the scalar strict bounded Gram for
all homogeneous states in the matched-angle chart.  Full normals
and the other testing charts remain separate requirements.
